\documentclass{article}
\usepackage{amsmath,amssymb}
\usepackage[vocab={continuedfractions}]{leanify}
\newtheorem{lemma}{Lemma}
\begin{document}

\section{Notation}

We write $a_n(x)$ for the $n$th partial quotient of $x$.
We write $\xi_n(x)$ for the $n$th complete quotient of $x$.
We write $c_n(x)$ for the $n$th convergent of $x$.
We write $q_n(x)$ for the denominator of the $n$th convergent of $x$.
We write $K_n(a)$ for the value of the continued fraction with partial quotients $a$ up to $n$.
For every real number $k$, let $e(k)\coloneqq 1$.

\section{Expansions}

\begin{lemma}\label{root}
$a_0(\sqrt2)=1$, $a_1(\sqrt2)=2$, $\xi_0(\sqrt2)=\sqrt2$ and $c_1(\sqrt2)=\frac32$.
\end{lemma}

\begin{lemma}\label{denominator}
Let $x$ be an irrational number. Then $q_0(x)=1$.
\end{lemma}

\begin{lemma}\label{finite}
$K_2(e)=\frac32$.
\end{lemma}

\section{Approximation}

\begin{lemma}\label{golden}
$\frac{1+\sqrt5}{2}$ is a quadratic irrational, and $\frac32$ is a best approximation of $\sqrt2$.
\end{lemma}

\begin{lemma}\label{lagrange}
Let $x$ be a real number. If $x$ is a quadratic irrational, then $x$ is badly approximable.
\end{lemma}

\begin{lemma}\label{liouville}
Every Liouville number is irrational.
\end{lemma}

\end{document}
